\documentclass[11pt,a4paper]{article}
\usepackage[top=3cm,bottom=3cm,left=2.5cm,right=2.5cm]{geometry}
\usepackage{graphicx}
\usepackage{booktabs}
\usepackage{longtable}
\usepackage{array}
\usepackage{pdflscape}
\usepackage{verbatim}
\usepackage{float}
\usepackage[hidelinks]{hyperref}

\newcommand{\todo}[1]{\textbf{[TODO: #1]}}
\newcommand{\code}[1]{\texttt{#1}}
\let\plainunderscore\_
\renewcommand{\_}{\plainunderscore\allowbreak}
\setlength{\emergencystretch}{3em}
\newcommand{\evidencelog}[1]{%
  \IfFileExists{#1}{{\footnotesize\verbatiminput{#1}}}{\todo{add \detokenize{#1}}}}
\newcommand{\evidencefigure}[3][0.9\linewidth]{%
  \begin{figure}[H]\centering
  \IfFileExists{#2}{\includegraphics[width=#1]{#2}}{\fbox{\todo{add \detokenize{#2}}}}
  \caption{#3}
  \end{figure}}

\title{\Large Big Data Analytics \\[0.5em]
\Huge Tests \\[0.5em]
\Large Real-time Financial Transaction Anomaly Detection}
\author{Team: philodendron white measure \\[0.5em]}

\begin{document}
\maketitle
\tableofcontents
\newpage

% ---------------------------------------------------------------------------------
\section{Scope and approach}

This document covers the testing deliverable of Milestone 2: functional tests of the
data acquisition and pre-processing modules, quality tests of the data transformation
processes, and the tabular documentation of every test case.

\subsection{Modules under test}

\begin{tabular}{@{}>{\raggedright\arraybackslash}p{0.38\linewidth}
  >{\raggedright\arraybackslash}p{0.58\linewidth}@{}}
\toprule
\textbf{Module} & \textbf{Role} \\
\midrule
\code{generator/prepare\_data.py} & Acquisition: batch/stream split of PaySim and
generation of synthetic login events. \\
\code{generator/stream\_generator.py} & Acquisition: step-wise replay of the stream to
Kafka with controlled fault injection. \\
\code{batch\_layer/clean\_stream\_data.py} & Pre-processing: de-duplication and imputation
of missing financial fields before storage in MinIO. \\
\code{batch\_layer/historical\_pipeline.py} & Transformation: validation, feature
engineering and aggregation of historical data in Spark. \\
\bottomrule
\end{tabular}

\subsection{Test categories}

\begin{description}
\item[Functional] The test checks that a module behaves as specified for a given input:
outputs, boundary cases, error handling, message order.
\item[Quality] The test checks a property of the data produced by a transformation:
completeness (no residual nulls), consistency (balance arithmetic, time attributes),
conservation (no row lost or invented), distribution (fault rates, scenario and device
shares) and reproducibility (determinism under the project seed). Quality tests carry
the \code{quality} pytest marker.
\end{description}

\subsection{Test levels}

\begin{description}
\item[Unit] \code{tests/unit}. Pure Python on small in-memory fixtures; Kafka clients are
replaced by fakes, no service is needed. Runs in CI on every pull request.
\item[Deployment check] \code{tests/integration} (\code{test\_minio\_setup},
\code{test\_kafka\_topics}, \code{test\_airflow\_dag}). Verifies that the Docker Compose
stack is provisioned as designed: buckets, topics and partitions, DAG registration and
schedule.
\item[Integration] \code{tests/integration/test\_stream\_to\_minio}. One end-to-end run of
the stream path: Kafka $\rightarrow$ Airflow DAG $\rightarrow$ cleaning $\rightarrow$ MinIO.
\end{description}

Deployment checks and the integration test share the \code{integration} marker and
appear with level \emph{integration} in Table~\ref{tab:test-cases}. They are skipped,
with the reason recorded, when their service is not reachable.

% ---------------------------------------------------------------------------------
\section{Test environment and commands}

\begin{tabular}{@{}p{0.30\linewidth}p{0.66\linewidth}@{}}
\toprule
\textbf{Component} & \textbf{Version} \\
\midrule
Kafka & \code{apache/kafka:4.3.1} \\
Airflow & \code{apache/airflow:2.5.0} \\
Cassandra & \code{cassandra:5.0} \\
\bottomrule
\end{tabular}

\subsection{Commands}

The full run (stack started with \code{docker compose up -d}):
\begin{verbatim}
uv run pytest tests -v -ra --test-doc docs/tests/test_cases \
    --html=docs/tests/evidence/pytest_report.html --self-contained-html \
    2>&1 | tee docs/tests/evidence/pytest_run.txt
\end{verbatim}
CI runs the same suite with \code{-m "not integration"}. The \code{--test-doc} option
generates the tables in Sections~\ref{sec:summary} and~\ref{sec:cases} from the test
docstrings and the outcomes of this run; they are not typed by hand.

% ---------------------------------------------------------------------------------
\section{Summary}
\label{sec:summary}

\begin{table}[H]
\centering
% Generated by pytest --test-doc on 2026-11-18 20:15 CET. Do not edit by hand.
\begin{tabular}{llrrr}
\toprule
\textbf{Level} & \textbf{Type} & \textbf{Tests} & \textbf{Cases} & \textbf{Passed} \\
\midrule
integration & functional & 8 & 8 & 8 \\
unit & functional & 35 & 45 & 45 \\
unit & quality & 11 & 11 & 11 \\
\midrule
Total &  & 54 & 64 & 64 \\
\bottomrule
\end{tabular}

\caption{Test functions and executed cases per level and type. A parametrised test is
one test with several cases.}
\label{tab:summary}
\end{table}

% ---------------------------------------------------------------------------------
\begin{landscape}
\section{Test cases}
\label{sec:cases}

One row per test function, in collection order. IDs \code{U-nn} are unit tests,
\code{I-nn} deployment checks and integration tests. The mapping from ID to pytest node
is in \code{test\_cases.md}.

% Generated by pytest --test-doc on 2026-11-18 19:39 CET. Do not edit by hand.
% Requires: booktabs, longtable, array.
\begingroup
\footnotesize
\setlength{\tabcolsep}{4pt}
\ifdefined\testdocwidth\else\newlength{\testdocwidth}\fi
\setlength{\testdocwidth}{\dimexpr\linewidth-16\tabcolsep}
\begin{longtable}{>{\raggedright\arraybackslash\hspace{0pt}}p{0.05\testdocwidth}>{\raggedright\arraybackslash\hspace{0pt}}p{0.10\testdocwidth}>{\raggedright\arraybackslash\hspace{0pt}}p{0.08\testdocwidth}>{\raggedright\arraybackslash\hspace{0pt}}p{0.08\testdocwidth}>{\raggedright\arraybackslash\hspace{0pt}}p{0.16\testdocwidth}>{\raggedright\arraybackslash\hspace{0pt}}p{0.21\testdocwidth}>{\raggedright\arraybackslash\hspace{0pt}}p{0.20\testdocwidth}>{\raggedright\arraybackslash\hspace{0pt}}p{0.12\testdocwidth}}
\caption{Test cases}\label{tab:test-cases}\\
\toprule
\textbf{ID} & \textbf{Module} & \textbf{Level} & \textbf{Type} & \textbf{Objective} & \textbf{Steps} & \textbf{Expected result} & \textbf{Actual result} \\
\midrule
\endfirsthead
\toprule
\textbf{ID} & \textbf{Module} & \textbf{Level} & \textbf{Type} & \textbf{Objective} & \textbf{Steps} & \textbf{Expected result} & \textbf{Actual result} \\
\midrule
\endhead
\midrule
\multicolumn{8}{r}{\emph{continued on next page}} \\
\endfoot
\bottomrule
\endlastfoot
I-01 & airflow\_\allowbreak{}dag & integration & functional & Verify Airflow parses the stream\_\allowbreak{}to\_\allowbreak{}minio DAG file without import errors. & 1. Run airflow dags list-import-errors -o json in the scheduler container. & No import error refers to the stream\_\allowbreak{}to\_\allowbreak{}minio DAG file. & passed \\
\addlinespace
I-02 & airflow\_\allowbreak{}dag & integration & functional & Verify the stream\_\allowbreak{}to\_\allowbreak{}minio DAG is registered in the scheduler. & 1. Run airflow dags list -o json in the scheduler container. & stream\_\allowbreak{}to\_\allowbreak{}minio is among the listed dag\_\allowbreak{}ids. & passed \\
\addlinespace
I-03 & airflow\_\allowbreak{}dag & integration & functional & Verify the DAG is scheduled as a one-minute micro-batch. & 1. Load the DAG with DagBag inside the scheduler container. \newline 2. Print schedule\_\allowbreak{}interval.total\_\allowbreak{}seconds(). & The schedule interval is 60 seconds. & passed \\
\addlinespace
I-04 & kafka\_\allowbreak{}topics & integration & functional & Verify kafka-init created the topics used by the pipeline. & 1. List the topics of the broker with the Kafka AdminClient. & raw-transactions, raw-logins and preprocessed-transactions exist. & passed \\
\addlinespace
I-05 & kafka\_\allowbreak{}topics & integration & functional & Verify each pipeline topic has the configured number of partitions. & 1. List the topics of the broker with the Kafka AdminClient. \newline 2. Count the partitions of the three pipeline topics. & Each topic has 3 partitions. & passed \\
\addlinespace
I-06 & minio\_\allowbreak{}setup & integration & functional & Verify run\_\allowbreak{}minio.py created every bucket the pipelines write to. & 1. List the MinIO buckets with boto3. & historical-data, clean-transactions-data, clean-logins-data and processed-data exist. & passed \\
\addlinespace
I-07 & minio\_\allowbreak{}setup & integration & functional & Verify the historical PaySim parquet is uploaded to MinIO. & 1. Request the metadata (head\_\allowbreak{}object) of paysim/historical\_\allowbreak{}paysim.parquet in the historical-data bucket. & - The object exists. \newline - Its size is greater than 0 bytes. & passed \\
\addlinespace
I-08 & stream\_\allowbreak{}to\_\allowbreak{}minio & integration & functional & Verify the stream path Kafka -\textgreater{} Airflow DAG -\textgreater{} cleaning -\textgreater{} MinIO end to end. & 1. Publish 10 uniquely marked transactions to raw-transactions, one of them with amount = null. \newline 2. Unpause and trigger the stream\_\allowbreak{}to\_\allowbreak{}minio DAG. \newline 3. Poll clean-transactions-data for up to 300 s for a new parquet that holds the batch. & - All 10 messages are delivered to Kafka and appear in one cleaned parquet. \newline - The financial columns of the batch contain no null. \newline - The blanked amount is imputed as 600.0. & passed \\
\addlinespace
U-01 & clean\_\allowbreak{}stream\_\allowbreak{}data & unit & functional & Verify identical stream messages collapse to one row. & 1. Build a batch with one message three times and a different one once. \newline 2. Call StreamConsumer.data\_\allowbreak{}cleaning. & - Two rows remain, one per account (C1, C2). \newline - No duplicated row is left. & passed \\
\addlinespace
U-02 & clean\_\allowbreak{}stream\_\allowbreak{}data & unit & functional & Verify a missing amount is imputed from the sender's balances. & 1. Build a message with amount = null (CASH\_\allowbreak{}OUT, CASH\_\allowbreak{}IN, PAYMENT cases). \newline 2. Clean it in a batch that also holds a complete reference row. & amount equals \textbar{}oldbalanceOrg - newbalanceOrig\textbar{} (600, 600 and 0). & 3 passed (3 cases) \\
\addlinespace
U-03 & clean\_\allowbreak{}stream\_\allowbreak{}data & unit & functional & Verify a missing old balance of an outflow transaction is imputed. & 1. Build a CASH\_\allowbreak{}OUT, PAYMENT or TRANSFER message (one case each) with oldbalanceOrg = null. \newline 2. Clean it in a batch that also holds a complete reference row. & oldbalanceOrg equals newbalanceOrig + amount (400 + 600 = 1000). & 3 passed (3 cases) \\
\addlinespace
U-04 & clean\_\allowbreak{}stream\_\allowbreak{}data & unit & functional & Verify a missing old balance of a CASH\_\allowbreak{}IN transaction is imputed. & 1. Build a CASH\_\allowbreak{}IN message with oldbalanceOrg = null. \newline 2. Clean it in a batch that also holds a complete reference row. & oldbalanceOrg equals newbalanceOrig - amount (1000 - 600 = 400). & passed \\
\addlinespace
U-05 & clean\_\allowbreak{}stream\_\allowbreak{}data & unit & functional & Verify a missing new balance of an outflow transaction is imputed. & 1. Build a CASH\_\allowbreak{}OUT, PAYMENT or TRANSFER message (one case each) with newbalanceOrig = null. \newline 2. Clean it in a batch that also holds a complete reference row. & newbalanceOrig equals oldbalanceOrg - amount (1000 - 600 = 400). & 3 passed (3 cases) \\
\addlinespace
U-06 & clean\_\allowbreak{}stream\_\allowbreak{}data & unit & functional & Verify a missing new balance of a CASH\_\allowbreak{}IN transaction is imputed. & 1. Build a CASH\_\allowbreak{}IN message with newbalanceOrig = null. \newline 2. Clean it in a batch that also holds a complete reference row. & newbalanceOrig equals oldbalanceOrg + amount (400 + 600 = 1000). & passed \\
\addlinespace
U-07 & clean\_\allowbreak{}stream\_\allowbreak{}data & unit & functional & Verify complete messages pass through cleaning unchanged. & 1. Build two complete messages (CASH\_\allowbreak{}OUT and CASH\_\allowbreak{}IN). \newline 2. Call data\_\allowbreak{}cleaning. & Output values and row order equal the input (dtypes aside). & passed \\
\addlinespace
U-08 & clean\_\allowbreak{}stream\_\allowbreak{}data & unit & functional & Verify login events are de-duplicated and otherwise left untouched. & 1. Build a batch with the same login event twice. \newline 2. Call data\_\allowbreak{}cleaning. & - One row remains. \newline - Its columns and values equal the input event. & passed \\
\addlinespace
U-09 & clean\_\allowbreak{}stream\_\allowbreak{}data & unit & functional & Verify no null remains in financial columns after single-field faults. & 1. Build 12 messages: amount, oldbalanceOrg and newbalanceOrig blanked once per type (CASH\_\allowbreak{}OUT, PAYMENT, TRANSFER, CASH\_\allowbreak{}IN). \newline 2. Call data\_\allowbreak{}cleaning. & - All 12 rows are kept. \newline - The three financial columns contain no null. & passed \\
\addlinespace
U-10 & clean\_\allowbreak{}stream\_\allowbreak{}data & unit & functional & Verify imputed rows satisfy the balance equation. & 1. Build six messages with one blanked financial field each (outflow types and CASH\_\allowbreak{}IN). \newline 2. Call data\_\allowbreak{}cleaning and compute old + sign * amount - new, with sign +1 for CASH\_\allowbreak{}IN and -1 otherwise. & The residual is 0 for every row. & passed \\
\addlinespace
U-11 & clean\_\allowbreak{}stream\_\allowbreak{}data & unit & functional & Verify cleaning works when a financial column is null in every row. & 1. Build two messages with the same financial column blanked (amount, oldbalanceOrg and newbalanceOrig cases). \newline 2. Call data\_\allowbreak{}cleaning. & - Both rows are kept. \newline - The column contains no null after cleaning. & 3 passed (3 cases) \\
\addlinespace
U-12 & historical\_\allowbreak{}pipeline & unit & functional & Verify validation rejects a frame that lacks a required column. & 1. Drop the amount column from the three-row sample frame. \newline 2. Call validate\_\allowbreak{}data on the reduced frame. & ValueError whose message starts with "Missing columns" is raised. & passed \\
\addlinespace
U-13 & historical\_\allowbreak{}pipeline & unit & functional & Verify engineered features are computed from the raw columns as specified. & 1. Call engineer\_\allowbreak{}features on the sample frame. \newline 2. Read the row of account C1 (TRANSFER of 100 from balance 500 to 400). & - amount\_\allowbreak{}log equals log(1 + 100) = 4.61512. \newline - hour\_\allowbreak{}of\_\allowbreak{}day is 0 and day is 1 for event\_\allowbreak{}time 2026-01-01 00:02:00. \newline - orig\_\allowbreak{}balance\_\allowbreak{}diff and dest\_\allowbreak{}balance\_\allowbreak{}diff both equal 100.0. \newline - amount\_\allowbreak{}to\_\allowbreak{}balance\_\allowbreak{}ratio equals 100 / 500 = 0.2. & passed \\
\addlinespace
U-14 & historical\_\allowbreak{}pipeline & unit & functional & Verify a zero originating balance yields a null ratio, not a division error. & 1. Set oldbalanceOrg to 0 in every row of the sample frame. \newline 2. Call engineer\_\allowbreak{}features and collect amount\_\allowbreak{}to\_\allowbreak{}balance\_\allowbreak{}ratio. & amount\_\allowbreak{}to\_\allowbreak{}balance\_\allowbreak{}ratio is null in every row (no infinity, no exception). & passed \\
\addlinespace
U-15 & historical\_\allowbreak{}pipeline & unit & functional & Verify per-step and per-type aggregates conserve counts and amounts. & 1. Call aggregate\_\allowbreak{}data on the sample frame (two rows in step 1, one in step 2). \newline 2. Read the step-1 row of hourly\_\allowbreak{}stats and the TRANSFER row of type\_\allowbreak{}stats. & - Step 1: transaction\_\allowbreak{}count 2, total\_\allowbreak{}amount 150.0, avg\_\allowbreak{}amount 75.0. \newline - TRANSFER: transaction\_\allowbreak{}count 2, total\_\allowbreak{}amount 300.0, avg\_\allowbreak{}amount 150.0. & passed \\
\addlinespace
U-16 & prepare\_\allowbreak{}data & unit & functional & Verify the batch/stream split loses and duplicates no PaySim row. & 1. Take the 12-row PaySim fixture with steps on both sides of step 500. \newline 2. Call prepare\_\allowbreak{}paysim(paysim\_\allowbreak{}df, 500). & - Batch holds only steps \textless{}= 500 and stream only steps \textgreater{} 500. \newline - Row counts add up to the input; the two index sets are disjoint. & passed \\
\addlinespace
U-17 & prepare\_\allowbreak{}data & unit & functional & Verify the boundary step 500 is assigned to the historical batch. & 1. Call prepare\_\allowbreak{}paysim(paysim\_\allowbreak{}df, 500); the fixture has a row at step 500. \newline 2. Collect the distinct step values of both outputs. & Step 500 occurs in the batch and not in the stream. & passed \\
\addlinespace
U-18 & prepare\_\allowbreak{}data & unit & functional & Verify the added time attributes are consistent with the PaySim step. & 1. Call prepare\_\allowbreak{}paysim(paysim\_\allowbreak{}df, 500). \newline 2. Compare event\_\allowbreak{}time and second\_\allowbreak{}offset with step in both outputs. & - event\_\allowbreak{}time equals (step - 1) * 3600 in every row; step 1 maps to 0. \newline - second\_\allowbreak{}offset equals 0 in every row. & passed \\
\addlinespace
U-19 & prepare\_\allowbreak{}data & unit & functional & Verify prepare\_\allowbreak{}paysim leaves the caller's DataFrame unchanged. & 1. Copy the PaySim fixture. \newline 2. Call prepare\_\allowbreak{}paysim on the original frame. & The original frame equals the copy (same columns and values). & passed \\
\addlinespace
U-20 & prepare\_\allowbreak{}data & unit & functional & Verify the split at extreme and intermediate split steps. & 1. Call prepare\_\allowbreak{}paysim with split\_\allowbreak{}step 0, 499 and 743 (one case each). \newline 2. Count the rows of batch and stream. & - Batch holds 0, 3 and 12 rows respectively. \newline - Stream holds the remaining rows of the 12-row fixture. & 3 passed (3 cases) \\
\addlinespace
U-21 & prepare\_\allowbreak{}data & unit & functional & Verify select\_\allowbreak{}accounts aggregates the stream per sending account. & 1. Split the fixture at step 500 and pass the stream to select\_\allowbreak{}accounts. \newline 2. Read the rows of accounts C\_\allowbreak{}B and C\_\allowbreak{}D. & - Output has account\_\allowbreak{}id, first\_\allowbreak{}transaction\_\allowbreak{}step and transaction\_\allowbreak{}count. \newline - C\_\allowbreak{}B: first step 503, 2 transactions; C\_\allowbreak{}D: first step 520, 2. \newline - There is exactly one row per distinct nameOrig. & passed \\
\addlinespace
U-22 & prepare\_\allowbreak{}data & unit & functional & Verify n\_\allowbreak{}accounts limits the selection to that many existing senders. & 1. Split the fixture at step 500. \newline 2. Call select\_\allowbreak{}accounts(stream, n\_\allowbreak{}accounts=3). & - Exactly 3 accounts are returned. \newline - Every returned account\_\allowbreak{}id is a nameOrig of the stream. & passed \\
\addlinespace
U-23 & prepare\_\allowbreak{}data & unit & functional & Verify every account receives one complete IEEE-CIS device profile. & 1. Select all stream accounts of the PaySim fixture. \newline 2. Call sample\_\allowbreak{}device\_\allowbreak{}profiles with the identity fixture, which mixes complete and incomplete rows. & - There is one output row per account and no device column is null. \newline - Every sampled profile exists among the complete identity rows. & passed \\
\addlinespace
U-24 & prepare\_\allowbreak{}data & unit & functional & Verify device identifiers are unique and sequential. & 1. Call sample\_\allowbreak{}device\_\allowbreak{}profiles for all stream accounts of the fixture. & device\_\allowbreak{}id runs D000001, D000002, ... in row order without gaps. & passed \\
\addlinespace
U-25 & prepare\_\allowbreak{}data & unit & functional & Verify sampled device types follow the reference distribution. & 1. Build 5000 synthetic accounts. \newline 2. Call sample\_\allowbreak{}device\_\allowbreak{}profiles with the identity fixture (non-null DeviceType: desktop 3/5, mobile 2/5). \newline 3. Compute the share of each sampled DeviceType. & Desktop share is 0.60 +/- 0.03 and mobile share is 0.40 +/- 0.03. & passed \\
\addlinespace
U-26 & prepare\_\allowbreak{}data & unit & functional & Verify identity data without any complete profile is rejected. & 1. Build identity data whose only row has null device columns. \newline 2. Call sample\_\allowbreak{}device\_\allowbreak{}profiles for one account. & ValueError is raised. & passed \\
\addlinespace
U-27 & prepare\_\allowbreak{}data & unit & functional & Verify activity patterns and login scenarios take documented values. & 1. Call assign\_\allowbreak{}login\_\allowbreak{}scenarios for 100 accounts. & - There is one output row per account. \newline - activity\_\allowbreak{}pattern is SHORT\_\allowbreak{}ACTIVITY or MIXED\_\allowbreak{}ACTIVITY and login\_\allowbreak{}scenario is a key of LOGIN\_\allowbreak{}SCENARIOS. \newline - The input frame gets no login\_\allowbreak{}scenario column. & passed \\
\addlinespace
U-28 & prepare\_\allowbreak{}data & unit & functional & Verify scenario and activity shares match the documented mix. & 1. Call assign\_\allowbreak{}login\_\allowbreak{}scenarios for 20 000 accounts. \newline 2. Compute the share of each activity pattern and login scenario. & - SHORT\_\allowbreak{}ACTIVITY share is 0.70 +/- 0.015. \newline - NORMAL\_\allowbreak{}LOGIN, SINGLE\_\allowbreak{}FAILURE, MULTIPLE\_\allowbreak{}FAILURES and MULTIPLE\_\allowbreak{}LOGINS shares are 0.65, 0.20, 0.05 and 0.10, each +/- 0.015. & passed \\
\addlinespace
U-29 & prepare\_\allowbreak{}data & unit & functional & Verify the schema and identifiers of generated login events. & 1. Call generate\_\allowbreak{}login\_\allowbreak{}events(accounts\_\allowbreak{}with\_\allowbreak{}devices, 500). & - Columns are the 12 documented ones, in the documented order. \newline - event\_\allowbreak{}id runs E0000001, E0000002, ... in row order. & passed \\
\addlinespace
U-30 & prepare\_\allowbreak{}data & unit & functional & Verify accounts without a pre-transaction window get no login events. & 1. Generate events for the fixture, which holds accounts with first transaction step 501 (C\_\allowbreak{}SKIP\_\allowbreak{}501) and 502 (C\_\allowbreak{}EDGE\_\allowbreak{}502). & - C\_\allowbreak{}SKIP\_\allowbreak{}501 has no events. \newline - C\_\allowbreak{}EDGE\_\allowbreak{}502 has events. & passed \\
\addlinespace
U-31 & prepare\_\allowbreak{}data & unit & functional & Verify each account's event types follow its login scenario. & 1. Generate events for the fixture (one account per scenario). \newline 2. Order the events of each eligible account by event\_\allowbreak{}time. & - Event types equal LOGIN\_\allowbreak{}SCENARIOS[login\_\allowbreak{}scenario], in order. \newline - event\_\allowbreak{}time is strictly increasing within an account. & passed \\
\addlinespace
U-32 & prepare\_\allowbreak{}data & unit & functional & Verify login events fall between the split and the first transaction. & 1. Generate events for the fixture. \newline 2. Compute each account's transaction start as (first\_\allowbreak{}transaction\_\allowbreak{}step - 1) * 3600. & - event\_\allowbreak{}time \textgreater{}= 500 * 3600 for every event. \newline - event\_\allowbreak{}time is lower than the transaction start of its account. & passed \\
\addlinespace
U-33 & prepare\_\allowbreak{}data & unit & functional & Verify SHORT\_\allowbreak{}ACTIVITY events stay within six hours of the transaction. & 1. Generate events for the fixture. \newline 2. Keep the events of SHORT\_\allowbreak{}ACTIVITY accounts. & event\_\allowbreak{}time \textgreater{}= transaction start - NEAR\_\allowbreak{}WINDOW\_\allowbreak{}HOURS * 3600 (6 h). & passed \\
\addlinespace
U-34 & prepare\_\allowbreak{}data & unit & functional & Verify gaps between consecutive events respect the scenario limits. & 1. Generate events for the fixture. \newline 2. Compute the gaps between consecutive events of each scenario account. & - NORMAL\_\allowbreak{}LOGIN has a single event, hence no gap. \newline - Gaps are 1 to 120 s for SINGLE\_\allowbreak{}FAILURE and 1 to 60 s for MULTIPLE\_\allowbreak{}FAILURES. \newline - Gaps are 1 s to 3 h for MULTIPLE\_\allowbreak{}LOGINS. & passed \\
\addlinespace
U-35 & prepare\_\allowbreak{}data & unit & functional & Verify step and second\_\allowbreak{}offset are consistent with event\_\allowbreak{}time. & 1. Generate events for the fixture. \newline 2. Recompute step and second\_\allowbreak{}offset from event\_\allowbreak{}time. & - step equals event\_\allowbreak{}time // 3600 + 1 for every event. \newline - second\_\allowbreak{}offset equals event\_\allowbreak{}time \% 3600 for every event. & passed \\
\addlinespace
U-36 & prepare\_\allowbreak{}data & unit & functional & Verify events carry the account's device profile under the new names. & 1. Generate events for the fixture. \newline 2. Compare each event with the row of its account. & - device\_\allowbreak{}id equals the account's device\_\allowbreak{}id. \newline - device\_\allowbreak{}type, operating\_\allowbreak{}system, browser, device\_\allowbreak{}info and screen\_\allowbreak{}resolution equal DeviceType, id\_\allowbreak{}30, id\_\allowbreak{}31, DeviceInfo and id\_\allowbreak{}33. & passed \\
\addlinespace
U-37 & prepare\_\allowbreak{}data & unit & functional & Verify login events are returned in chronological order. & 1. Generate events for the fixture. & event\_\allowbreak{}time is non-decreasing over the output rows. & passed \\
\addlinespace
U-38 & prepare\_\allowbreak{}data & unit & functional & Verify an account with a one-hour window gets its events in that hour. & 1. Generate events for the fixture. \newline 2. Select the events of C\_\allowbreak{}EDGE\_\allowbreak{}502 (MULTIPLE\_\allowbreak{}LOGINS, first step 502). & - Two events exist. \newline - Both lie in [500 * 3600, 501 * 3600). & passed \\
\addlinespace
U-39 & prepare\_\allowbreak{}data & unit & functional & Verify save\_\allowbreak{}data writes the three parquet outputs with all rows. & 1. Redirect PROCESSED\_\allowbreak{}DIR and RAW\_\allowbreak{}STREAM\_\allowbreak{}DIR to a temporary directory. \newline 2. Build batch, stream and login events and call save\_\allowbreak{}data. \newline 3. Read the three parquet files back. & - paysim\_\allowbreak{}batch.parquet, paysim\_\allowbreak{}stream.parquet and login\_\allowbreak{}events.parquet exist. \newline - Each file has as many rows as the frame it was written from. & passed \\
\addlinespace
U-40 & prepare\_\allowbreak{}data & unit & functional & Verify login-event generation is reproducible under the project seed. & 1. Seed NumPy with RANDOM\_\allowbreak{}SEED and run split, account selection, device sampling, scenario assignment and event generation. \newline 2. Repeat the run with the same seed. & - The first result is not empty. \newline - Both results are identical frames. & passed \\
\addlinespace
U-41 & stream\_\allowbreak{}generator & unit & functional & Verify stream inputs are accepted only as parquet paths. & 1. Call StreamProducer.\_\allowbreak{}check\_\allowbreak{}data\_\allowbreak{}path with a path ending in x.parquet. \newline 2. Call it with a path ending in x.csv. & - The parquet path is returned with suffix .parquet. \newline - The csv path raises ValueError. & passed \\
\addlinespace
U-42 & stream\_\allowbreak{}generator & unit & functional & Verify only rows of the current simulation step are selected. & 1. Build a producer at step 502 over rows of steps 501 and 502. \newline 2. Call \_\allowbreak{}take\_\allowbreak{}next\_\allowbreak{}stream. & Two rows are returned and both have step 502. & passed \\
\addlinespace
U-43 & stream\_\allowbreak{}generator & unit & functional & Verify fault injection neither loses nor invents transactions. & 1. Seed random with 0 and build 20 000 transactions of one step. \newline 2. Call \_\allowbreak{}data\_\allowbreak{}corruption and read the rows stored as delayed. & - Emitted rows = input rows - delayed rows + duplicates. \newline - Delayed rows are stored for the step and are absent from the emitted originals. \newline - The input frame is unchanged. & passed \\
\addlinespace
U-44 & stream\_\allowbreak{}generator & unit & functional & Verify injected fault rates match the documented values. & 1. Seed random with 1 and corrupt 20 000 transactions of one step. \newline 2. Count rows with a blanked field, duplicated rows and delayed rows. & - Blanked share is 2 \% +/- 0.5 pp, with at most one blanked field per row. \newline - Duplicated share is 0.5 \% +/- 0.2 pp. \newline - Delayed share is 1 \% +/- 0.3 pp. & passed \\
\addlinespace
U-45 & stream\_\allowbreak{}generator & unit & functional & Verify a fault-free replay sends every row once, in step order. & 1. Disable faults by patching random.random to return 1.0. \newline 2. Run produce() over 3 rows of step 501, 2 rows of step 502 and 1 login. \newline 3. Read the messages recorded by the fake Kafka producer. & - 5 transactions (steps 501, 501, 501, 502, 502) and 1 login are sent. \newline - Messages carry all input columns. \newline - The replay stops at the empty step 503. & passed \\
\addlinespace
U-46 & stream\_\allowbreak{}generator & unit & functional & Verify a delayed transaction is re-emitted in the following step. & 1. Script random.random so that only the first row of step 501 is delayed. \newline 2. Run produce() over steps 501 and 502. & - All 5 transactions are sent, none is lost. \newline - Step order is 501, 501, 502, 502, 501; the last message is account C0. & passed \\
\addlinespace
\end{longtable}
\endgroup

\end{landscape}

% ---------------------------------------------------------------------------------
\appendix
\section{Evidence}

\subsection{pytest run log}
\evidencelog{evidence/pytest_run.txt}

\subsection{Screenshots}
\evidencefigure{evidence/minio/buckets_present.png}{MinIO console: the four buckets
created by the batch layer.}
\evidencefigure{evidence/minio/clean_transaction_data.png}{MinIO console: cleaned
parquet in \code{clean-transactions-data} written by the end-to-end test.}
\evidencefigure{evidence/airflow/airflow_dag_run.png}{Airflow: successful run of the
\code{stream\_to\_minio} DAG triggered by the end-to-end test.}
\evidencefigure{evidence/kafka/kafka_topics.png}{Kafka UI: topics and partitions.}

\end{document}
